% Shrnutí okruhu Multiagentní systémy — jedna otázka = jedna strana.
% Kondenzát z 01-multiagentni-systemy.tex; slouží k opakování těsně před zkouškou.
% Struktura strany: znění otázky → osnova odpovědi → fakta v blocích → pasti.
\documentclass[10pt,a4paper]{article}
\usepackage[utf8]{inputenc}
\usepackage[T1]{fontenc}
\usepackage{lmodern}          % vektorové fonty — nutné pro expansion v microtype
\usepackage[czech]{babel}
\usepackage[margin=1.25cm,top=1.15cm,bottom=1.25cm]{geometry}
\usepackage{amsmath,amssymb}
\usepackage{parskip}
\usepackage{enumitem}
\setlist{nosep,leftmargin=1.1em,labelsep=0.35em,topsep=1pt}
\usepackage{multicol}
\setlength{\columnsep}{5mm}
\usepackage{booktabs}
\usepackage{array}
\usepackage{xcolor}
\definecolor{linkblue}{RGB}{20,60,150}
\definecolor{defbg}{RGB}{240,244,252}
\definecolor{defframe}{RGB}{100,130,190}
\definecolor{headcol}{RGB}{25,70,140}
\definecolor{markcol}{RGB}{176,84,0}
\usepackage[most]{tcolorbox}
\usepackage{microtype}
\usepackage[colorlinks=true,linkcolor=linkblue,urlcolor=linkblue]{hyperref}
\usepackage{fancyhdr}

\pagestyle{fancy}
\fancyhf{}
\renewcommand{\headrulewidth}{0.4pt}
\fancyhead[L]{\footnotesize\textsc{Multiagentní systémy --- shrnutí ke SZZ}}
\fancyhead[R]{\footnotesize\thepage}

% Nadpis otázky: číslo + doslovné znění věty oficiálního okruhu.
\newtcolorbox{qbox}[1]{colback=defbg,colframe=defframe,boxrule=0.7pt,arc=1.2mm,
  left=2.5mm,right=2.5mm,top=1.2mm,bottom=1.2mm,#1}
\newcommand{\qpage}[2]{%
  \begin{qbox}{}
  \textbf{\large Otázka #1.}\quad\textbf{#2}
  \end{qbox}
  \vspace{-0.3em}}
\newcommand{\hh}[1]{\par\smallskip\noindent{\bfseries\color{headcol}#1}\par\nobreak\vspace{-0.2em}}
\newcommand{\ph}[1]{\par\smallskip\noindent{\bfseries\color{markcol}#1}\par\nobreak\vspace{-0.2em}}
\newcommand{\ok}[1]{{\bfseries\color{markcol}#1)}\,}   % krok osnovy odpovědi
\newcommand{\tm}[1]{\textbf{#1}}
\newcommand{\en}[1]{\emph{#1}}
\newcommand{\ra}{$\rightarrow$}

\begin{document}
\raggedcolumns   % nestahovat/neroztahovat sloupce svisle

%=====================================================================
% Strana 1 — rozcestník
%=====================================================================
\thispagestyle{empty}
\begin{center}
{\LARGE\bfseries Multiagentní systémy}\\[0.35em]
{\large Shrnutí okruhu ke státní závěrečné zkoušce --- jedna otázka na jednu stranu}\\[0.5em]
{\small SZZ Umělá inteligence, zaměření Inteligentní agenti, MFF UK \quad$\cdot$\quad srpen 2026}
\end{center}

\begin{qbox}{}
\small\raggedright
\textbf{Oficiální znění okruhu} (zdroj pravdy, 2025/2026):
{\footnotesize\url{https://www.mff.cuni.cz/cs/studenti/bc-a-mgr-studium/studijni-plany/2025-2026/informatika/mgr/umela-inteligence}}
\end{qbox}

\small
Sedm vět okruhu, sedm stran. Podrobný výklad v~\path{01-multiagentni-systemy.pdf}
(anglicky \path{ENG/01-multiagent-systems.pdf}); zdrojové slidy \path{materialy/}
(NAIL106, Neruda, verze 2025).

\begin{enumerate}[leftmargin=2em,itemsep=1pt]
\item Architektura autonomního agenta; mechanismus výběru akcí. \hfill \textit{str.~2}
\item Metody pro řízení agentů; symbolické a~konekcionistické reaktivní plánování, hybridní
  přístupy. \hfill \textit{str.~3}
\item Problém hledání cesty. \hfill \textit{str.~4}
\item Komunikace a~znalosti v~multiagentních systémech, ontologie, řečové akty, FIPA-ACL,
  protokoly. \hfill \textit{str.~5}
\item Distribuované řešení problémů, kooperace, Nashova ekvilibria, Paretova efektivita,
  alokace zdrojů, aukce. \hfill \textit{str.~6}
\item Metody pro učení agentů; zpětnovazební učení. \hfill \textit{str.~7}
\item Metodologie návrhu, jazyky a~prostředí multiagentních systémů. \hfill \textit{str.~8}
\end{enumerate}

\begin{multicols}{2}
\small

\hh{Čtyři osy, kolem kterých se otázky točí}
\begin{itemize}
\item \tm{Reaktivita vs.\ proaktivita} --- definice inteligentního agenta, strategie odhodlání
  (bold/cautious), parametry $\gamma/\phi$ u~Maes, vrstvy hybridních architektur.
\item \tm{Individuální vs.\ kolektivní racionalita} --- sobec vede na $(D,D)$, tragédii obecní
  pastviny a~taktické hlasování; odtud mechanism design (Vickrey, VCG): navrhnout
  \emph{pravidla}, za nichž se pravda vyplatí.
\item \tm{Symbolické vs.\ subsymbolické} --- FOL, STRIPS, ontologie, ACL \emph{versus} subsumpce,
  aktivační sítě, NN v~RL, dnes LLM; praxe je hybridní.
\item \tm{Teorie vs.\ praxe} --- nekonstruktivní věty (optimální agent, Nash, Arrow) versus
  algoritmy (A*, CBS, Q-learning, CNET, JADE). Složitost: NE je PPAD-úplné, WDP, Slater
  i~optimální MAPF NP-těžké.
\end{itemize}

\hh{Co se chce spočítat nebo nakreslit u~tabule}
\begin{itemize}
\item Smíšené NE hry $2\times2$ (indiferentní \emph{protihráč}) --- typicky Battle of sexes.
\item Aukce s~oceněními 80/60/30 pro všechny čtyři typy; důkaz pravdomluvnosti Vickreyho.
\item Jeden krok Q-learningu / SARSA s~konkrétními čísly; cliff walking.
\item A* na mřížce: $f=g+h$, pořadí expanzí, manhattanská vs.\ oktilová heuristika.
\item Borda / Condorcet z~profilu preferencí; ukázat cyklus $A{>}B{>}C{>}A$.
\item Subsumpční pravidla Steelsova průzkumníka a~jejich priorita.
\end{itemize}

\columnbreak

\hh{Co v~slidech NENÍ (a přesto to okruh jmenuje)}
Hledání cesty a~MAPF (ot.~3), metodologie Gaia/Prometheus/Tropos a~platforma FIPA
(ot.~7), VCG a~Zeuthen, epistemická logika, DisCSP. Čerpáno z~Wooldridge a~z~kurzů
NAIL071/NAIL069; ve velké verzi je to označené.
\emph{Nové ve verzi 2025:} kapitola \tm{LLM agenti} (str.~8).

\ph{Nejčastější chyby}
\begin{itemize}
\item Zaměnit \en{desire} a~\en{intention} --- záměr je \emph{závazek}: přetrvává, vede
  means-ends reasoning a~omezuje další deliberaci.
\item Tvrdit, že efektem \texttt{inform} je, že příjemce uvěří. U~Cohena--Perraulta je efekt
  \uv{$H$ věří, že $S$ věří~$\varphi$}; u~FIPA je $B_j\varphi$ jen \emph{zamýšlený} (RE), nikoli
  vynutitelný.
\item Zaměnit přípustnou a~konzistentní heuristiku (konzistence je silnější a~dává expanzi
  každého vrcholu nejvýš jednou).
\item Říct, že $(D,D)$ je špatné, protože není NE. Je to NE --- a~zároveň jediné
  \emph{ne}-Pareto-optimální pole matice.
\item Připsat FPSB a~holandské aukci dominantní strategii (nemají ji).
\item Zaměnit on-policy a~off-policy: SARSA se učí hodnotu politiky, kterou hraje;
  Q-learning hodnotu optimální politiky bez ohledu na sběr dat.
\item Plést směr subsumpční priority --- \emph{nižší} (reaktivnější) vrstva přebíjí vyšší.
\end{itemize}

\end{multicols}

\newpage

%=====================================================================
\qpage{1}{Architektura autonomního agenta; mechanismus výběru akcí}
%=====================================================================
\begin{multicols}{2}
\small

\hh{Osnova odpovědi}
\ok{1} co je agent a~MAS \ok{2} tři vlastnosti $+$ intencionální postoj \ok{3} vlastnosti
prostředí \ok{4} abstraktní architektura $\langle E,e_0,\tau\rangle$ a~$Ag$ \ok{5} dva
speciální případy \ok{6} jak agentovi zadat cíl (utilita, predikát) \ok{7} mechanismus výběru
akcí a~jeho rodiny \ra{} most k~otázce~2.

\hh{Agent a~MAS}
\tm{Agent} (Wooldridge--Jennings) $=$ počítačový systém zasazený do prostředí, schopný
\emph{autonomní akce} vedoucí ke splnění \emph{delegovaných} cílů.
\tm{Franklin--Graesser}: systém, který je \emph{součástí} prostředí, vnímá je, jedná v~něm
\emph{v~čase} a~sleduje vlastní agendu tak, aby ovlivnil, co bude vnímat příště.
\tm{AIMA}: cokoli, co vnímá senzory a~jedná efektory (nejširší, a~proto nejméně užitečná).
\tm{MAS} $=$ systém interagujících agentů, kteří \emph{nemusí} mít sdílený cíl \ra{} konflikty,
vyjednávání. Motivace 90.~let: ubikvitnost, konektivita, složitost, delegace.
Vymezení: \emph{objekt} nemá kontrolu nad svým chováním (\uv{objekty to dělají zadarmo, agenti
proto, že chtějí}); v~\emph{distribuovaném systému} je koordinace daná předem, agenti ji řeší
až za~běhu; \emph{expertní systém} není zasazen do prostředí a~sám nejedná.

\hh{Tři vlastnosti inteligentního agenta}
\tm{reaktivita} (včasná reakce na změnu) $\cdot$ \tm{proaktivita} (vlastní iniciativa a~cíle)
$\cdot$ \tm{sociálnost} (interakce). Těžké je jejich \emph{vyvážení}: čistý reaktivec nikam
nedojde, čistý proaktivec ignoruje realitu.

\hh{Intencionální postoj (Dennett)}
Tři postoje: \tm{fyzikální} (zákony) $\cdot$ \tm{návrhový} (účel komponent) $\cdot$
\tm{intencionální} (přesvědčení a~přání). Intencionální postoj je \emph{abstrakce pro zvládnutí
složitosti}: legitimní, pokud pomáhá predikovat; \emph{není} tvrzením o~vědomí.

\hh{Vlastnosti prostředí}
plně / \tm{částečně pozorovatelné} $\cdot$ statické / \tm{dynamické} $\cdot$ deterministické /
\tm{nedeterministické} $\cdot$ diskrétní / \tm{spojité}. Reálná prostředí bývají v~té těžší
variantě a~z~toho přímo plyne obtížnost návrhu agenta.

\hh{Abstraktní architektura}
$E$ stavy, $A$ akce, \tm{běh} $r=e_0a_0e_1a_1\dots$; $R^A$ běhy končící akcí, $R^E$ stavem.
\tm{Přechodová funkce} $\tau:R^A\to 2^E$ --- doménou \emph{celá historie}, obrazem
\emph{množina} (nedeterminismus), $\tau(r)=\emptyset$ značí konec.
Prostředí $\mathit{Env}=\langle E,e_0,\tau\rangle$, agent $Ag:R^E\to A$, systém
$\langle Ag,\mathit{Env}\rangle$ s~množinou běhů $\mathcal{R}(Ag,\mathit{Env})$.
\tm{Funkční ekvivalence} $=$ shodné množiny běhů (vůči danému prostředí, resp.\ všem).

\hh{Dva speciální případy}
\tm{Čistě reaktivní agent} $Ag:E\to A$ --- rozhoduje jen podle aktuálního stavu, bez historie;
\emph{striktně slabší}.
\tm{Agent s~vnitřním stavem}: $\mathit{see}:E\to\mathit{Per}$,
$\mathit{next}:I\times\mathit{Per}\to I$, $\mathit{action}:I\to A$, smyčka
\emph{vnímej\,\ra\,aktualizuj\,\ra\,jednej}. Oba lze převést na standardního agenta $R^E\to A$;
opačně u~čistě reaktivního nikoli.

\hh{Jak agentovi říct, co má dělat}
\tm{Utilita stavu} $u:E\to\mathbb{R}$ --- neocení \emph{jak} se stavu dosáhlo;
\tm{utilita běhu} $u:R\to\mathbb{R}$ --- vyjádří průběh, ale zadává se neintuitivně.
\tm{Optimální agent} maximalizuje očekávanou utilitu
$\sum_r P(r\mid Ag,\mathit{Env})\,u(r)$ --- definice \emph{nekonstruktivní} a~předpokládající
neomezené zdroje (realističtější je \en{bounded optimality}).
Příklad \tm{Tileworld}: $u(r)=N_f/N_a$ (splněné/všechny úkoly).
\tm{Predikátová specifikace}: $\langle\mathit{Env},\Psi\rangle$, $\Psi:R\to\{0,1\}$,
$\mathcal{R}_\Psi$ úspěšné běhy. \tm{Pesimista} uspěje jen tehdy, jsou-li úspěšné
\emph{všechny} běhy, \tm{optimista} stačí-li jeden, \tm{realista} počítá pravděpodobnost.
Dva typy úloh: \tm{dosažení} (\en{achievement}, dostaň se do $G$) a~\tm{udržení}
(\en{maintenance}, nikdy nespadni do $B$).

\hh{Mechanismus výběru akcí --- jádro otázky}
Procedura, kterou agent v~každém kroku rozhodne, \emph{kterou akci provede}, z~vjemů
a~vnitřního stavu. \tm{Architektura} $=$ softwarová realizace takového mechanismu, tj.\ rozklad
na moduly a~jejich interakce (Maes 1991). Rodiny --- osnova otázky~2:
\begin{itemize}
\item \tm{symbolická/deliberativní} --- dedukce a~plánování: deduktivní agent, STRIPS, BDI, PRS;
\item \tm{reaktivní symbolická} --- pravidla \emph{situace\,\ra\,akce} s~pevnou prioritou:
  subsumpce (Brooks);
\item \tm{reaktivní konekcionistická} --- šíření aktivace, vyhraje nejaktivnější modul:
  agent network architecture (Maes);
\item \tm{hybridní} --- vrstvy různé abstrakce $+$ arbitráž: TouringMachines, InteRRaP;
\item \tm{koaliční} --- soutěž koalic procesů o~\uv{vědomí}: IDA;
\item \tm{naučená} --- maximalizace $Q$ nebo stochastická politika: Q-learning, actor-critic;
\item \tm{generativní} --- LLM jako politika i~plánovač (nově ve slidech 2025, str.~8).
\end{itemize}

\end{multicols}
\newpage

%=====================================================================
\qpage{2}{Metody pro řízení agentů; symbolické a~konekcionistické reaktivní plánování,
hybridní přístupy}
%=====================================================================

\vspace{-0.2em}
\begin{center}\footnotesize
\begin{tabular}{@{}lllll@{}}
\toprule
\textbf{rodina} & \textbf{reprezentace} & \textbf{arbitráž} & \textbf{hlavní slabina} &
\textbf{zástupci} \\
\midrule
deliberativní & symbolická, model světa & dedukce, plánovač &
  transdukce, čas na výpočet & deduktivní agent, STRIPS, BDI, PRS \\
reaktivní & žádná / subsymbolická & priorita, aktivace &
  žádný výhled do budoucna & subsumpce, Maes \\
hybridní & obojí, po vrstvách & mediátor či průchod vrstev &
  složitý návrh, rozhraní & TouringMachines, InteRRaP \\
koaliční & blackboard $+$ koalice & soutěž koalic &
  ad hoc volby, ladění & IDA \\
\bottomrule
\end{tabular}
\end{center}
\vspace{-0.3em}

\begin{multicols}{2}
\small

\hh{Osnova odpovědi}
\ok{1} deduktivní agent a~jeho dvě fatální slabiny \ok{2} BDI $+$ STRIPS $+$ odhodlání
\ok{3} PRS \ok{4} východiska reaktivního přístupu \ok{5} subsumpce (Steels) \ok{6} Maes
\ok{7} hybridní architektury \ok{8} IDA \ok{9} shrnout čtyřmi typy arbitráže.

\hh{Deduktivní agent}
Stav $=$ databáze formulí $\mathit{DB}\in 2^L$; výběr akce $=$ důkaz
$\mathit{DB}\vdash_\rho \mathit{Do}(a)$. Dvě fáze: najdi \emph{dokazatelnou} akci, jinak akci
\emph{konzistentní} s~$\mathit{DB}$, jinak nic. Elegantní a~verifikovatelný, ale nepraktický:
\tm{problém transdukce} (dostat svět do symbolů včas) a~\tm{vypočitatelná racionalita}
(než důkaz doběhne, svět se změní).

\hh{BDI --- praktické uvažování}
\tm{Beliefs} subjektivní a~omylné $\cdot$ \tm{desires} smějí být nekonzistentní $\cdot$
\tm{intentions} jsou \emph{závazky}: přetrvávají, vedou means-ends reasoning, omezují další
deliberaci, musí být konzistentní s~přesvědčeními.
Dvě fáze: \tm{deliberace} ($\mathit{brf}$\,\ra\,$\mathit{options}$\,\ra\,$\mathit{filter}$,
tj.\ \emph{co chci}) a~\tm{means-ends reasoning} (plánování, \emph{jak toho dosáhnout}).
Ve smyčce navíc $\mathit{reconsider}$ (vyplatí se přehodnotit?) a~$\mathit{sound}$
(platí plán ještě?).

\hh{STRIPS}
Deskriptor akce $\langle P_a,D_a,A_a\rangle$ (předpoklady, delete, add), problém
$\langle B_0,O,G\rangle$, efekt $B_i=(B_{i-1}\setminus D_{a_i})\cup A_{a_i}$.
Plán \tm{přípustný} $\Leftrightarrow$ $B_{i-1}\models P_{a_i}$ pro všechna $i$; \tm{korektní}
$\Leftrightarrow$ navíc $B_n\models G$. Umět projít svět kostek.

\hh{Odhodlání (\en{commitment})}
Kdy záměr opustit: \tm{blind} (až po dosažení) $\cdot$ \tm{single-minded} (dokud věřím, že je
dosažitelný) $\cdot$ \tm{open-minded} (dokud ho chci).
\tm{Bold} přehodnocuje až po dokončení plánu, \tm{cautious} po každém kroku; efektivita
$=$ dosažené/všechny záměry. Pravidlo: \emph{pomalu se měnící prostředí\,\ra\,bold, rychle
se měnící\,\ra\,cautious.}

\hh{PRS}
První implementace BDI. Plán $=$ \tm{Goal} $+$ \tm{Context} $+$ \tm{Body}; místo plánování se
z~\emph{knihovny plánů} bere plán s~platným kontextem; zásobník záměrů; deliberace metaplány
nebo utilitou. Potomci: AgentSpeak/Jason, JAM, JACK, Jadex.

\hh{Reaktivní východiska (Brooks)}
\tm{situatedness} (agent je v~prostředí, ne v~modelu) $\cdot$ \tm{embodiment} (tělo a~senzory)
$\cdot$ \tm{emergence} (inteligence z~interakce jednoduchých chování).
Teze: \uv{svět je svůj nejlepší model}; inteligentní chování lze získat bez symbolické
reprezentace i~bez explicitního uvažování.

\hh{Symbolické reaktivní plánování: subsumpce}
Množina jednoduchých \tm{chování} \emph{situace\,\ra\,akce}, běžících \emph{paralelně},
s~\tm{pevnou prioritou} $\prec$ ($b_1\prec b_2$ znamená \uv{$b_1$ inhibuje $b_2$}, tedy výčet
od nejvyšší priority). Provede se nejvýše prioritní chování, které \uv{hoří}; nehoří-li nic,
neděje se nic. Rychlé (konstantní složitost, lze v~HW), robustní, ale s~desítkami chování
těžko laditelné.
\emph{Steelsův průzkumník Marsu} --- základna vysílá navigační signál, roboti nekomunikují
přímo:
\begin{itemize}\itemsep0pt
\item R1 překážka \ra{} změň směr; R2 nesu vzorek a~jsem na základně \ra{} polož;
  R3 nesu a~nejsem \ra{} po gradientu nahoru; R4 vidím vzorek \ra{} seber; R5 jinak náhodná chůze.
\item Druhá iterace přidává \tm{stigmergii} přes \uv{drobky}: R7 nesu vzorek \ra{} \emph{polož
  2~drobky} a~jdi po gradientu; R8 cítím drobek \ra{} \emph{seber~1} a~jdi po gradientu dolů.
  Priorita $R1\prec R6\prec R7\prec R4\prec R8\prec R5$.
\end{itemize}

\hh{Konekcionistické reaktivní plánování: Maes}
\tm{Agent network architecture}: moduly s~\emph{předpoklady}, \emph{add/delete listem}
a~\emph{prahem aktivace}, propojené \tm{následnickými}, \tm{předchůdcovskými}
a~\tm{konfliktními} spoji. Aktivace přitéká ze \emph{stavu světa} a~z~\emph{cílů}, inhibice
z~\emph{chráněných cílů}; provede se \emph{proveditelný} modul s~nejvyšší aktivací nad prahem
(pak se vynuluje). Parametry $\gamma$ (šíření z~cílů) a~$\phi$ (ze světa) ladí proaktivitu
proti reaktivitě.

\hh{Omezení reaktivních architektur}
Bez modelu světa nelze uvažovat o~budoucnosti; jen krátkodobý a~lokální pohled; emergence není
inženýrská metodika (chování se ladí experimentálně); mnoho vrstev a~jejich interakcí neškáluje.

\hh{Hybridní architektury}
\tm{Horizontální}: vrstvy paralelně, každá vidí vjemy a~navrhuje akci, rozhoduje
\emph{mediátor} --- až $m^n$ kombinací, mediátor je úzké hrdlo.
\tm{Vertikální}: vrstvy sériově, \en{one-pass} nebo \en{two-pass}, $n-1$ rozhraní ---
jednodušší, ale křehké (selže vrstva \ra{} selže agent).
\tm{TouringMachines} --- horizontální, vrstvy modelling / planning / reactive $+$ řídicí
pravidla. \tm{InteRRaP} --- vertikální two-pass, vrstvy kooperativní plánování / lokální
plánování / behaviorální, každá s~vlastní KB (vjemy zdola nahoru, akce shora dolů).
\tm{Stanley} (vítěz DARPA Grand Challenge 2005) jako reálné vrstvené nasazení.

\hh{IDA a~globální pracovní prostor}
Franklin $+$ Baarsova \en{global workspace theory}: mysl jako MAS jednoduchých nevědomých
procesů komunikujících přes \emph{blackboard} (asociativní pole); \tm{koalice} procesů soutěží
o~vstup do \uv{vědomí}, vítězná se vykoná; hierarchie kontextů. Nejsložitější architektura,
mnoho ad hoc rozhodnutí a~obtížné ladění.

\ph{Čtyři typy arbitráže --- tím odpověď uzavřít}
\tm{pevná priorita} (subsumpce) $\cdot$ \tm{úroveň aktivace} (Maes) $\cdot$
\tm{mediační funkce} (vrstvené) $\cdot$ \tm{soutěž koalic} (IDA).

\end{multicols}
\newpage

%=====================================================================
\qpage{3}{Problém hledání cesty}
%=====================================================================
\begin{multicols}{2}
\small

\hh{Osnova odpovědi}
\ok{1} formulace, diskretizace prostoru \ok{2} neinformované vs.\ informované prohledávání
\ok{3} A* a~vlastnosti heuristiky \ok{4} varianty a~meze \ok{5} dynamické prostředí: LPA*,
D*~Lite \ok{6} MAPF --- konflikty, cílové funkce, složitost \ok{7} algoritmy MAPF.

\hh{Formulace a~zařazení}
Prohledávání stavového prostoru grafu $G=(V,E)$ s~ohodnocením $w:E\to\mathbb{R}^+$, řešením je
nejlevnější cesta. V~MAS jde o~nejjednodušší \emph{means-ends reasoning}: agent má model
prostředí a~hledá posloupnost akcí k~cíli.
\emph{Ve slidech NAIL106 kapitola není --- okruh ji ale jmenuje; zdroj je AIMA a~NAIL071.}

\hh{Jak vzniknou vrcholy grafu}
\tm{Mřížka} (4- nebo 8-okolí) $\cdot$ \tm{viditelnostní graf} (rohy překážek --- dá optimální
cestu v~mnohoúhelníkovém světě) $\cdot$ \tm{navigační síť} (navmesh, konvexní plochy) $\cdot$
\tm{stavová mřížka} (poloha $+$ orientace $+$ rychlost, pro neholonomní roboty) $\cdot$
\tm{vzorkovací metody} pro spojitý a~vysokodimenzionální prostor (PRM, RRT, RRT*) $\cdot$
\tm{potenciálová pole} (velmi rychlá, ale uvíznou v~lokálním minimu).

\hh{Škála prohledávacích algoritmů}
BFS (optimální jen při jednotkových cenách) $\cdot$ \tm{Dijkstra/UCS} ($h\equiv0$) $\cdot$
\en{greedy best-first} (jen $h$ --- rychlé, \emph{neoptimální}) $\cdot$ \tm{A*} (kombinuje
obojí) $\cdot$ obousměrné prohledávání. Nadstavby: \tm{JPS} (odstraní symetrie v~mřížce),
\tm{any-angle} Theta* (cesta nemusí ležet na hranách mřížky), \tm{HPA*} (hierarchie clusterů).

\hh{A*}
$f(n)=g(n)+h(n)$: $g$ známá cena do~$n$, $h$ odhad zbytku. Prioritní fronta \en{open},
množina \en{closed}; expanduje se vrchol s~nejmenším~$f$.
\begin{itemize}
\item \tm{Přípustnost}: $h(n)\le h^*(n)$ --- nikdy nepřestřelí \ra{} A* je \emph{optimální}
  (u~grafové verze je třeba dovolit reexpanzi).
\item \tm{Konzistence} (monotonie): $h(n)\le w(n,m)+h(m)$ \ra{} $f$ neklesá podél cesty
  a~\emph{každý vrchol se expanduje nejvýš jednou}. Konzistence $+$ $h(\text{cíl})=0$
  \ra{} přípustnost; obráceně to neplatí.
\item $h\equiv 0$ dá \tm{Dijkstru}; $f=g+\varepsilon h$ (\tm{weighted A*}) je rychlejší, ale
  jen $\varepsilon$-suboptimální; \tm{IDA*} šetří paměť za cenu opakovaných průchodů.
\item Heuristiky v~mřížce: manhattanská (4-okolí), oktilová (8-okolí), euklidovská (libovolný
  úhel). Silnější přípustná heuristika \ra{} méně expanzí.
\end{itemize}

\hh{Dynamické prostředí}
Agent nezná mapu předem nebo se mapa mění; přeplánovat od nuly je drahé.
\tm{LPA*}: u~každého vrcholu drží $g$ a~\tm{rhs} (odhad z~předchůdců); vrchol je
\emph{lokálně konzistentní}, když $g=\mathit{rhs}$; přepočítávají se jen nekonzistentní vrcholy
zasažené změnou.
\tm{D*~Lite}: LPA* prohledávající \emph{od cíle k~agentovi} (proto pohyb agenta nevyžaduje
přepočet celého prostoru), s~inkrementální opravou po každém vjemu; \tm{freespace assumption}
--- neznámé buňky se považují za volné.

\hh{MAPF --- formulace}
$k$ agentů, každý se startem a~cílem, hledá se \emph{bezkonfliktní} sada cest.
Konflikty: \tm{vrcholový} (dva agenti ve stejném vrcholu v~čase~$t$) a~\tm{hranový}/swap
(prohodí se po téže hraně). Cílové funkce: \tm{sum-of-costs} (součet dob) vs.\ \tm{makespan}
(čas posledního) --- optimum jedné nemusí být optimem druhé.
Optimální MAPF je \tm{NP-těžký}.

\hh{Algoritmy MAPF}
\begin{itemize}
\item \tm{Prioritizované plánování / Cooperative A*}: agenti se plánují postupně podle priority
  do \emph{rezervační tabulky} v~prostoru-čase. Velmi rychlé, ale \emph{neúplné} a~neoptimální
  --- špatné pořadí priorit může udělat řešitelnou instanci neřešitelnou.
\item \tm{CBS}: \emph{nízká úroveň} $=$ A* v~prostoru-čase pro jednoho agenta se seznamem
  omezení; \emph{vysoká úroveň} $=$ strom omezení --- najde se konflikt a~vytvoří dvě větve,
  každá zakazuje jednomu z~agentů daný vrchol/hranu v~daném čase.
  \emph{Optimální a~úplný}; dvě úrovně se vyhnou exponenciálnímu sdruženému prostoru
  $|A|^k$.
\item Škálovatelné varianty: ECBS (bounded suboptimal), SIPP (safe interval), LNS
  (large neighborhood search).
\end{itemize}

\ph{Pasti}
Přípustnost $\ne$ konzistence. Weighted A* není optimální. Prioritizované plánování není
úplné. \uv{Optimální MAPF} vždy doplnit \emph{vzhledem ke které cílové funkci}.

\hh{Vazby}
Cesty se plánují jako means-ends reasoning (ot.~2); koordinace tras je speciální případ
distribuovaného řešení problémů a~DisCSP (ot.~5); místo plánování lze politiku
\emph{naučit} (ot.~6).

\end{multicols}
\newpage

%=====================================================================
\qpage{4}{Komunikace a~znalosti v~multiagentních systémech, ontologie, řečové akty,
FIPA-ACL, protokoly}
%=====================================================================
\begin{multicols}{2}
\small

\hh{Osnova odpovědi}
\ok{1} proč komunikace $\ne$ volání metody \ok{2} řečové akty: Austin, Searle
\ok{3} Cohen--Perrault --- sémantika ve stylu STRIPS \ok{4} KQML/KIF a~FIPA-ACL, FP a~RE
\ok{5} interakční protokoly \ok{6} ontologie a~jazyky \ok{7} epistemická logika.

\hh{Proč je komunikace agentů jiná}
Zpráva není volání metody: agent je autonomní, \emph{příjemce není povinen vyhovět}.
Komunikace je proto \tm{akce} --- mění stav světa (přesvědčení příjemce) a~lze ji plánovat
stejně jako fyzickou akci.

\hh{Řečové akty}
\tm{Austin} 1962: vyslovením se \emph{koná} (\uv{prohlašuji vás za muže a~ženu}).
\tm{Lokuce} (co bylo řečeno) / \tm{ilokuce} (jaký akt to byl) / \tm{perlokuce} (jaký efekt to
mělo).
\tm{Searle}: podmínky zdařilosti (normální vstup/výstup, propoziční obsah, přípravné,
upřímnost, esenciální) a~pět kategorií --- \tm{asertiva} (informování), \tm{direktiva}
(žádost o~akci), \tm{komisiva} (slib), \tm{expresiva} (postoj), \tm{deklarace} (mění stav
světa). Zpráva $=$ \emph{performativní sloveso} $+$ \emph{propoziční obsah}.

\hh{Cohen \& Perrault (1979)}
Řečový akt jako operátor \emph{STRIPS} (předpoklady / efekt) s~modálními operátory
\emph{believe}, \emph{cando}, \emph{want}.
\emph{Request}$(S,H,A)$: pre --- $S$ věří, že $H$ umí~$A$, \emph{a} $S$ věří, že $H$ věří, že
umí~$A$ (plus $S$ věří, že to chce); efekt --- $H$ věří, že $S$ věří, že $S$ chce~$A$.
\emph{Inform}$(S,H,\varphi)$: pre --- $S$ věří~$\varphi$; efekt --- \emph{$H$ věří, že $S$
věří~$\varphi$} (nikoli, že $H$ uvěří~$\varphi$!).

\hh{KQML a~FIPA-ACL}
\tm{KQML} $=$ \emph{obálka} (ilokuční část): performativum $+$ \texttt{:sender},
\texttt{:receiver}, \texttt{:content}, \texttt{:language}, \texttt{:ontology},
\texttt{:reply-with}, \texttt{:in-reply-to}; obsah v~\tm{KIF} (FOL v~lispovské syntaxi).
Zprostředkování přes \emph{facilitátory} (\texttt{recruit-one}, \texttt{broker-all},
\texttt{advertise}). Kritika: chybí \texttt{commit}, nejasná sémantika, roztažitá množina
performativ.
\tm{FIPA-ACL} (standard, zjednodušení KQML): \texttt{performative},
\texttt{sender}/\texttt{receiver}, \texttt{content}, \texttt{language}, \texttt{ontology},
\texttt{protocol}, \texttt{conversation-id}, \texttt{reply-by}; \tm{22 performativ}
(\texttt{inform}, \texttt{request}, \texttt{cfp}, \texttt{propose}, \texttt{accept-} /
\texttt{reject-proposal}, \texttt{agree}, \texttt{refuse}, \texttt{failure},
\texttt{not-understood}, \texttt{subscribe}, \texttt{proxy}, \dots).
Sémantika v~jazyce \tm{SL}: \tm{FP} (\en{feasibility precondition}) a~\tm{RE} (\en{rational
effect}). Pro \texttt{inform} je RE${}=B_j\varphi$ --- efekt \emph{zamýšlený}, ale
nevynutitelný (Cohenův--Perraultův efekt je slabší). Hlavní teoretický problém:
\tm{neverifikovatelnost} --- zvenčí nelze zkontrolovat, že agent opravdu věří tomu, co tvrdí.

\hh{Interakční protokoly}
Předepsané posloupnosti zpráv, tedy konečný automat konverzace:
\tm{FIPA-Request} (\texttt{request}\,\ra\,\texttt{agree}/\texttt{refuse}\,\ra\,
\texttt{inform}/\texttt{failure}), \tm{FIPA-Query}, \tm{Contract Net} a~jeho iterovaná varianta
(\texttt{cfp}\,\ra\,\texttt{propose}/\texttt{refuse}\,\ra\,\texttt{accept}/\texttt{reject}\,
\ra\,\texttt{inform}), \tm{Subscribe}, \tm{Brokering/Recruiting}, aukční protokoly
(English, Dutch). Protokol se uvádí polem \texttt{protocol}, vlákno drží
\texttt{conversation-id}.

\hh{Ontologie}
\tm{Gruber}: ontologie $=$ \emph{explicitní specifikace konceptualizace} --- formalizovaný
popis části reality, na kterém se agenti shodnou.
Stavební kameny: \tm{třídy}, \tm{instance}, \tm{vlastnosti}, \tm{relace} (zejména tranzitivní
\emph{is-a}) a~\tm{fakta}. Platí \emph{ontologie $+$ fakta $=$ znalostní báze}.
Škála formalizace: řízený slovník \ra{} taxonomie \ra{} tezaurus \ra{} rámce \ra{} logická
omezení. Dělení na \tm{aplikační}, \tm{doménové} a~\tm{horní} (\en{upper}, např.\ SUMO).

\hh{Ontologické jazyky}
\tm{RDF}: trojice \emph{subjekt--predikát--objekt} (graf), dotazy SPARQL; \tm{RDFS} přidává
třídy, hierarchii a~doménu/rozsah vlastností.
\tm{OWL}: \emph{deskripční logika} $=$ rozhodnutelný fragment FOL; \tm{open world assumption}
(co není řečeno, není nepravda); \tm{T-Box} (terminologie) vs.\ \tm{A-Box} (fakta).
Profily OWL~1: Lite ($\mathcal{SHIF}$), DL ($\mathcal{SHOIN}$), Full (nerozhodnutelný);
OWL~2 $=\mathcal{SROIQ}$ $+$ profily EL/QL/RL.

\hh{Epistemická logika}
$K_i\varphi$ $=$ \uv{agent~$i$ ví~$\varphi$}, sémantika \emph{možných světů} (relace
dosažitelnosti). \tm{Znalost} $=$ \tm{S5} (obsahuje axiom~T: $K_i\varphi\to\varphi$ --- znát
znamená mít pravdu), \tm{přesvědčení} $=$ \tm{KD45} (bez~T, lze věřit nepravdě; D konzistence,
4/5 introspekce). \tm{Společná znalost} $C_G\varphi$ \emph{nevznikne} nespolehlivou komunikací
--- \tm{problém koordinovaného útoku}.

\end{multicols}
\newpage

%=====================================================================
\qpage{5}{Distribuované řešení problémů, kooperace, Nashova ekvilibria, Paretova efektivita,
alokace zdrojů, aukce}
%=====================================================================

\vspace{-0.2em}
\begin{center}\footnotesize
\begin{tabular}{@{}lllll@{}}
\toprule
\textbf{aukce} & \textbf{forma} & \textbf{cena} & \textbf{dominantní strategie} &
\textbf{80/60/30} \\
\midrule
anglická   & open cry, rostoucí & první & přihazuj $\varepsilon$ až po své ocenění
  \emph{(slabě dom.)} & 61 \\
holandská  & open cry, klesající & první & \emph{neexistuje} --- kdy zastavit závisí na
  soupeřích & 80 \\
FPSB       & sealed bid & první & \emph{neexistuje} --- o~kolik podstřelit závisí na soupeřích
  & 80 \\
Vickrey    & sealed bid & druhá & nabídni \emph{pravdivé} ocenění \emph{(slabě dom.)} & 60 \\
\bottomrule
\end{tabular}\\[0.15em]
{\footnotesize Strategická ekvivalence: anglická $\leftrightarrow$ Vickrey,
holandská $\leftrightarrow$ FPSB. Sloupec 80/60/30 platí bez jakékoli informace o~soupeřích.}
\end{center}
\vspace{-0.3em}

\begin{multicols}{2}
\small

\hh{Osnova odpovědi}
\ok{1} koherence vs.\ koordinace, CDPS/PPS/MAS \ok{2} sdílení úloh a~výsledků: CNET, BBS,
DisCSP \ok{3} sobecký agent a~hra v~normální formě \ok{4} NE, Nashova věta, PPAD
\ok{5} Pareto a~sociální blahobyt \ok{6} klasické hry \ok{7} sociální volba a~Arrow
\ok{8} aukce, VCG, vyjednávání.

\hh{Koherence, koordinace, tři paradigmata}
\tm{Koherence} $=$ jak si vede systém \emph{jako celek}; \tm{koordinace} $=$ minimalizace režie
na synchronizaci a~konfliktů. Koordinace je prostředek, koherence cíl.
\tm{CDPS}: agenti \emph{benevolentní} se \emph{sdíleným cílem} (benevolence zásadně
zjednodušuje návrh). \tm{PPS}: homogenní procesory, jeden algoritmus, centrální řízení.
\tm{MAS}: vlastní cíle, konflikty, vyjednávání.
Tři fáze CDPS: \tm{dekompozice}\,\ra\,\tm{řešení podproblémů}\,\ra\,\tm{syntéza}.

\hh{Sdílení úloh a~výsledků}
\tm{Sdílení úloh} --- kdo co udělá, se určí \emph{dohodou} (CNET). \tm{Sdílení výsledků} ---
proaktivní nebo reaktivní; přínosy: \tm{jistota} (nezávislé potvrzení), \tm{úplnost} (skládání
lokálních pohledů), \tm{přesnost}, \tm{včasnost}.
\tm{CNET} (Smith \& Davis 1977): \tm{recognition}\,\ra\,\tm{announcement} (cfp se specifikací
a~omezeními)\,\ra\,\tm{bidding}\,\ra\,\tm{awarding/expediting}. Role \emph{manažer}
a~\emph{dodavatel} jsou dynamické, subkontrakty se hierarchicky vnořují; standardizováno jako
FIPA-Contract-Net.
\tm{BBS}: \emph{tabule} $=$ sdílená struktura s~hierarchií úrovní abstrakce (+ mutex),
\emph{experti} přispívají tím, co umí, \emph{arbitr} vybírá, kdo bude pracovat. Volná vazba
mezi experty; tabule je úzké hrdlo a~vynucuje společnou architekturu. (Užito i~v~IDA.)
\tm{DisCSP/DCOP}: proměnné a~omezení rozdělené mezi agenty, \emph{nikdo nevidí celý problém};
\tm{ABT} posílá \texttt{ok?} a~\texttt{nogood} podle pevného pořadí agentů; DCOP maximalizuje
\emph{sociální blahobyt} (ADOPT, DPOP, max-sum).

\hh{Sobecký agent a~hra v~normální formě}
Sobec maximalizuje \emph{svoji} očekávanou utilitu za předpokladu, že ostatní dělají totéž ---
společná utilita se tím \emph{nemaximalizuje}, ale strategie je robustní.
\tm{Hra} $\langle N,(A_i),(u_i)\rangle$, matice výplat; \tm{ryzí} vs.\ \tm{smíšené} strategie;
\tm{dominance} (lepší při \emph{každé} volbě ostatních), \tm{iterované odstraňování
dominovaných}.

\hh{Nash a~Pareto}
\tm{NE} $=$ profil vzájemně nejlepších odpovědí; \emph{nikdo se nemá důvod jednostranně
odchýlit}. Naivní hledání prochází $m^n$ profilů.
\tm{Nashova věta}: každá konečná hra má NE ve \emph{smíšených} strategiích (nekonstruktivní);
hledání je \tm{PPAD-úplné} (2006) --- \uv{o~něco lepší než NP-úplné}.
\tm{Pareto-optimalita}: nelze zlepšit jednomu, aniž se jinému přitíží.
\tm{Sociální blahobyt} $=\sum_i u_i$; maximum sw je Pareto-optimální, obráceně to neplatí.
\emph{Smíšené NE pro $2\times2$}: pravděpodobnosti volím tak, aby byl \emph{protihráč}
indiferentní.

\hh{Klasické hry (umět matice zpaměti)}
\tm{Vězňovo dilema} $3/3$, $0/4$, $4/0$, $1/1$, obecně $T>R>P>S$: $D$ dominantní,
$(D,D)$ jediné NE a~\emph{jediné ne-Pareto-optimální} pole, $(C,C)$ maximalizuje sw.
Analogie \emph{tragedy of the commons}. \tm{Iterované PD}: \emph{stín budoucnosti} umožňuje
spolupráci (známý konečný počet kol ji zpětnou indukcí zase odstraní); Axelrodův turnaj
a~\tm{TFT} --- slušná, odvetná, odpouštějící, nezávistivá, srozumitelná.
\tm{Koordinační hra} (Battle of sexes) $1/2$, $2/1$, jinak $0/0$: dvě ryzí NE $(B,B)$, $(F,F)$;
smíšené NE hraje \emph{svoji preferovanou} akci s~pravd.\ $2/3$ a~dá užitek jen $2/3$ ---
tedy \emph{méně} než kterékoli ryzí NE.
\tm{Anti-koordinační hra} (Chicken / Hawk--Dove) $5/5$, $1/6$, $6/1$, $0/0$: ryzí NE jsou
$(C,D)$ a~$(D,C)$, sociální blahobyt ale maximalizuje $(C,C)$.
\tm{Matching pennies}: žádné ryzí NE, jen smíšené.

\hh{Sociální volba}
Pravidla: plurality, plurality s~eliminací (IRV), sekvenční párové, \tm{Borda}
($N-1,\dots,0$), Bordovy kombinace (Black, Baldwin, Nanson), \tm{Slater} (minimalizuje počet
obrácených hran turnaje, NP-těžké).
\tm{Condorcetův paradox}: $A{>}B{>}C$, $B{>}C{>}A$, $C{>}A{>}B$ --- individuální preference
tranzitivní, kolektivní cyklická (analogie kámen-nůžky-papír).
Kritéria a~kdo je splňuje: \tm{Pareto} --- plurality a~Borda ano, sekvenční párové ne;
\tm{Condorcetův vítěz} --- sekvenční párové a~Bordovy kombinace ano, Borda samotná ne;
\tm{IIA} --- \emph{nesplňuje žádné} z~nich.
\tm{Arrowova věta}: nad 3$+$ alternativami je jediné pravidlo splňující Pareto $+$ IIA
(při univerzalitě) \emph{diktatura}. \tm{Gibbard--Satterthwaite}: každý rozumný systém lze
\emph{manipulovat} taktickým hlasováním.

\hh{Aukce a~návrh mechanismů}
\emph{Efektivní} aukce přidělí zdroj tomu, kdo ho oceňuje nejvíc. Dimenze: private / common /
correlated value; open cry / sealed bid; first / second price; počet kol.
Umět \emph{důkaz pravdomluvnosti Vickreyho} (porovnat nadhodnocení a~podhodnocení oproti
pravdě --- ani jedno nikdy nezlepší výsledek).
\tm{Věta o~ekvivalenci výnosů}: očekávaný výnos všech čtyř aukcí je stejný --- neplatí při
averzi k~riziku, u~common values (\emph{prokletí vítěze}) a~při koluzi.
Další formáty: \emph{kombinatorické} (nabídky na množiny; \tm{WDP} je NP-těžký),
\emph{all-pay} (platí všichni --- model lobbingu), \emph{tichá}, \emph{amsterdamská}
(anglická, pak od dvou zájemců holandská).
\tm{Mechanism design} $=$ obrácená teorie her: navrhnout pravidla tak, aby žádoucí chování bylo
v~ekvilibriu; \emph{princip odhalení} (stačí zkoumat pravdomluvné mechanismy).
\tm{VCG}: alokace maximalizující sw, platba $=$ \emph{externalita}, kterou vítěz způsobí
ostatním \ra{} pravdomluvnost.
\tm{Vyjednávání}: monotónní protokol ústupků $+$ \tm{Zeuthenova strategie} --- ustupuje ten,
kdo riskuje méně.

\end{multicols}
\newpage

%=====================================================================
\qpage{6}{Metody pro učení agentů; zpětnovazební učení}
%=====================================================================
\begin{multicols}{2}
\small

\hh{Osnova odpovědi}
\ok{1} proč se agenti učí a~co se učí \ok{2} MDP, Bellman, value/policy iteration
\ok{3} Q-learning vs.\ SARSA \ok{4} průzkum vs.\ využití \ok{5} politikové metody a~hluboké RL
\ok{6} MARL: Markovská hra, algoritmy, potíže, CTDE.

\hh{Proč a~co}
Návrhář nezná prostředí dopředu, prostředí se mění, optimální chování je příliš složité na
ruční naprogramování.
Paradigmata: \tm{s~učitelem} $\cdot$ \tm{bez učitele} $\cdot$ \tm{zpětnovazební} --- v~MAS
nejběžnější, protože zpětná vazba bývá jen skalární odměna. Dále \tm{imitační učení}
a~\tm{inverzní RL} (odvoď odměnu z~chování experta), \tm{evoluční} metody (LCS, neuroevoluce).
Agent se učí \emph{model prostředí}, \emph{model protivníka}, \emph{politiku} nebo utilitu;
\tm{individuálně} nebo \tm{sociálně} (napodobování, sdílení zkušenosti).

\hh{Markovský rozhodovací proces}
$\mathit{MDP}=\langle S,A,T,R,\gamma\rangle$ s~\tm{Markovskou vlastností} (budoucnost závisí
jen na aktuálním stavu, ne na historii); $T:S\times A\times S\to[0,1]$,
$0\le\gamma<1$ váží současnost proti budoucnosti.
Politika $\pi:S\to A$ (nebo stochastická), výnos $G_t=\sum_k\gamma^k r_{t+k}$,
$V^\pi(s)=\mathbb{E}[G_t\mid s]$, $Q^\pi(s,a)$.
\tm{Bellmanova rovnice optimality}
$V^*(s)=\max_a\sum_{s'}T(s,a,s')\bigl[R+\gamma V^*(s')\bigr]$.
\tm{Value iteration} (Bellmanův operátor je kontrakce \ra{} konverguje),
\tm{policy iteration} (evaluace $+$ zlepšení, konečně mnoho iterací).
\emph{Známý model \ra{} plánování; neznámý \ra{} učení.}

\hh{Q-learning}
$$Q(s,a)\leftarrow Q(s,a)+\alpha\bigl[r+\gamma\max_{a'}Q(s',a')-Q(s,a)\bigr]$$
\tm{off-policy} (učí optimální $Q$ bez ohledu na chování při sběru dat), model-free,
tabulkový; konverguje při dostatečném průzkumu a~klesajícím~$\alpha$.
\emph{Umět spočítat jeden krok s~konkrétními čísly.}

\hh{SARSA}
$Q(s,a)\leftarrow Q(s,a)+\alpha[r+\gamma Q(s',a')-Q(s,a)]$ --- \tm{on-policy}: učí hodnotu
politiky, kterou skutečně hraje, a~proto je \emph{opatrnější}. V~úloze \en{cliff walking} volí
cestu dál od útesu, zatímco Q-learning jde po hraně (a při $\varepsilon$-greedy sběru dat tam
občas spadne).

\hh{Průzkum vs.\ využití}
$\varepsilon$-greedy $\cdot$ softmax/Boltzmann $\cdot$ optimistická inicializace $\cdot$ UCB.
Bez dostatečného průzkumu neplatí konvergenční záruky.

\hh{Metody založené na politice}
\tm{Policy gradient / REINFORCE}: parametrizovaná politika $\pi_\theta$,
$\nabla_\theta J=\mathbb{E}[\nabla_\theta\ln\pi_\theta(a\mid s)\,G_t]$ --- zvyš pravděpodobnost
akcí, které vedly k~vysokému výnosu; $G_t$ z~Monte Carlo epizod. Vhodné pro spojité akce.
\tm{Actor-critic}: \emph{herec} $=$ politika, \emph{kritik} $=$ hodnotová funkce, používá
\tm{výhodu} $A=Q-V$ (nižší rozptyl). \tm{A3C} $=$ asynchronní advantage actor-critic (více
pracovníků aktualizuje gradienty asynchronně).
\tm{DQN}: $Q$ jako neuronová síť $+$ \emph{replay buffer} (dekoreluje data) $+$ \emph{target
network} (stabilizuje cíl); rekurentní varianta pro částečnou pozorovatelnost.

\hh{Multiagentní RL}
\tm{Markovská hra} (\en{stochastic game}) $=$ MDP pro $N$ agentů: sdružený akční prostor
$A_1\times\dots\times A_N$, individuální odměny $R_i$; hodnota politiky agenta závisí na
politikách ostatních. Pojmy z~otázky~5 se přenášejí: nejlepší odpověď, NE, Pareto-optimalita.
\tm{Minimax-Q} (hry s~nulovým součtem --- v~každé aktualizaci se řeší LP pro minimaxovou
hodnotu), \tm{Nash-Q} (obecný součet, potřebuje NE podhry).
Potíže: \tm{nestacionarita} (ostatní se učí, prostředí se z~pohledu agenta mění),
\emph{exponenciální} sdružený akční prostor, částečná pozorovatelnost, velký rozptyl odhadů,
\tm{přiřazení zásluh}, oscilace mezi několika NE.
Řešení: \tm{CTDE} (centralizované učení, decentralizované vykonávání --- MADDPG, QMIX),
\emph{parameter sharing}, \tm{self-play} a~AlphaZero; \tm{MADRL} je dnešní standard.

\hh{Kam který algoritmus patří}
\begin{center}\footnotesize
\begin{tabular}{@{}llll@{}}
\toprule
& \textbf{model} & \textbf{politika} & \textbf{co drží} \\
\midrule
value / policy iteration & známý & --- & $V$, resp.\ $\pi$ \\
Monte Carlo & ne & on & $V$, $Q$ \\
SARSA & ne & on & $Q$ \\
Q-learning, DQN & ne & off & $Q$ \\
REINFORCE & ne & on & $\pi_\theta$ \\
actor-critic, A3C & ne & on & $\pi_\theta$ $+$ $V$ \\
\bottomrule
\end{tabular}
\end{center}

\hh{Dvě věci navíc}
\tm{Shaping odměny}: přidání potenciálové funkce $F(s,s')=\gamma\Phi(s')-\Phi(s)$
\emph{nemění} optimální politiku --- libovolné jiné doplnění odměny ji změnit může.
\tm{Smrtící trojice} (\en{deadly triad}): aproximace funkce $+$ bootstrapping $+$ off-policy
učení dohromady mohou divergovat; DQN to tlumí replay bufferem a~target sítí.

\ph{Pasti}
On-policy vs.\ off-policy (viz cliff walking). Value iteration je \emph{plánování}, ne učení.
Konvergence Q-learningu platí jen v~tabulkovém případě. V~MARL už \uv{optimální politika}
sama o~sobě nedává smysl --- jen vůči politikám ostatních.

\hh{Vazby}
Naučená $Q$-funkce nebo politika \emph{je} mechanismus výběru akcí (ot.~1); RL je alternativa
k~plánování (ot.~2) i~k~hledání cesty (ot.~3); MARL stojí na ekvilibriích (ot.~5).

\end{multicols}
\newpage

%=====================================================================
\qpage{7}{Metodologie návrhu, jazyky a~prostředí multiagentních systémů}
%=====================================================================
\begin{multicols}{2}
\small

\hh{Osnova odpovědi}
\ok{1} proč nestačí objektové metodologie \ok{2} Gaia \ok{3} Prometheus a~Tropos
\ok{4} standard FIPA a~architektura platformy \ok{5} jazyky a~prostředí \ok{6} LLM agenti
\ok{7} jak volit nástroj.

\hh{Proč nestačí objektové metodologie}
Agent je \tm{autonomní} (nelze mu zavolat metodu a~čekat vykonání), má \tm{vlastní cíle}
a~vztahy mezi agenty jsou \tm{organizační} (role, závazky, protokoly), ne strukturální
(dědičnost, kompozice). \tm{AOSE} proto pracuje s~pojmy role, organizace, interakce, cíl.

\hh{Gaia}
Metodologie pro \emph{uzavřené, kooperativní} systémy se statickou organizací.
\tm{Analýza}: \emph{model rolí} --- \en{permissions} (co role smí číst a~měnit),
\en{responsibilities} dělené na \tm{liveness} (co se má stát; výrazy s~$\cdot$, $|$, $\omega$)
a~\tm{safety} (invarianty), \en{activities}, \en{protocols}; plus \emph{model interakcí}.
\tm{Návrh}: model \emph{agentů} (typy agentů $\leftrightarrow$ role a~jejich instance), model
\emph{služeb}, model \emph{známostí} (\en{acquaintance}, kdo s~kým komunikuje).
Předpoklady: benevolentní agenti, žádné konflikty cílů, organizace se za běhu nemění.

\hh{Prometheus a~Tropos}
\tm{Prometheus} --- cíleno na BDI a~implementaci: systémová specifikace (aktéři, scénáře,
\emph{funkcionality}) \ra{} architektonický návrh (\emph{typy agentů}, protokoly) \ra{}
detailní návrh (\emph{capabilities}, plány, události, datové struktury); nástroj PDT.
\tm{Tropos} --- staví na požadavkovém inženýrství (i*, Eric Yu): \emph{aktéři}, \emph{cíle}
a~\tm{softgoals} (nefunkční požadavky), \emph{závislosti} mezi aktéry. Jediná z~trojice, která
pokrývá i~\emph{rané požadavky} (proč systém vzniká): early/late requirements \ra{}
architectural \ra{} detailed design.

\hh{Standard FIPA a~architektura platformy}
FIPA (1996, dnes IEEE) standardizuje ACL, protokoly, ontologické služby a~\emph{platformu}:
\begin{itemize}
\item \tm{AMS} --- \uv{bílé stránky}: povinná evidence agentů, \tm{AID} (jména), životní cyklus
  (create, suspend, resume, destroy);
\item \tm{DF} --- \uv{žluté stránky}: agenti registrují \emph{služby} a~vyhledávají
  poskytovatele;
\item \tm{MTS/ACC} --- doručování ACL zpráv uvnitř platformy i~mezi platformami.
\end{itemize}

\columnbreak

\hh{Jazyky a~prostředí}
\begin{itemize}
\item \tm{AOP} (Shoham 1993) --- historické východisko: stav agenta popsán přesvědčeními,
  závazky a~schopnostmi, program $=$ množina \emph{pravidel závazků}; první jazyk
  \tm{AGENT-0}. Jinou cestou jde \tm{Concurrent MetateM} (agent $=$ přímo spustitelná
  temporální specifikace).
\item \tm{JADE} --- Java, plně FIPA: třídy \texttt{Agent}, \texttt{Behaviour}
  (\texttt{Simple}/\texttt{Cyclic}/\texttt{FSM}), \texttt{ACLMessage}; kontejnery a~distribuce,
  hotové protokoly (Contract Net), nástroje Sniffer, Dummy, Introspector. Agent je vlákno,
  chování se plánují kooperativně --- \texttt{action()} nesmí blokovat.
\item \tm{Jason / AgentSpeak(L)} --- deklarativní \emph{BDI} jazyk, potomek PRS: plán
  \texttt{spouštěč : kontext <- tělo} (\texttt{+!has(owner,beer) : available(\dots) <- \dots}),
  přesvědčení jako logické literály, cíle \texttt{!} (dosáhni) a~\texttt{?} (zjisti);
  interpret běží BDI smyčku.
\item \tm{NetLogo} --- agentové \emph{simulace} (ABM): \emph{turtles}, \emph{patches},
  \emph{links}; jednoduché reaktivní agenty a~emergentní chování (Wolf--Sheep, Ants).
\item \tm{SPADE} --- Python, komunikace přes XMPP, \en{behaviours} jako v~JADE.
\item Další: JACK, Jadex, \tm{JaCaMo} (Jason $+$ CArtAgO pro prostředí a~artefakty $+$ Moise
  pro organizaci), MASON, GAMA, Mesa.
\end{itemize}

\hh{LLM agenti (nové ve slidech 2025)}
Velký jazykový model se používá jako:
\emph{politika} (mechanismus výběru akcí --- prompt obsahuje stav, výstupem je akce),
\emph{plánovač} (prohledávání prostoru plánů, dekompozice cíle),
\emph{kritik} v~actor-critic schématu, a~jako rozhraní pro \emph{volání nástrojů}
(tool use) a~pro \en{human-in-the-loop} učení.
Typické nasazení jsou \emph{agentové simulace}: LLM agenti jako uzly grafu sociálních
interakcí (šíření zpráv sítí) nebo jako \emph{syntetický panel} populace podle
sociologických průzkumů (odhad výsledků voleb).
Vazba na okruh: jde o~další rodinu mechanismů výběru akcí (ot.~1) --- bez explicitního modelu
světa, neverifikovatelná, ale s~velmi širokou doménovou znalostí.

\hh{Jak volit}
Simulace emergentního chování \ra{} NetLogo, MASON; průmyslový distribuovaný systém s~ACL
\ra{} JADE; agenti s~explicitními cíli a~plány \ra{} Jason, Jadex; rychlý prototyp v~Pythonu
\ra{} SPADE, Mesa.

\end{multicols}

\end{document}
